\section{The conductor and the normalization defect}
\label{sec:conductor}

We next bound the length of the adjugate subscheme by the singularities
of the plane curve.  For a reduced projective curve \(C\), let
\(\nu:\widetilde C\to C\) be its normalization and set
\begin{equation}
 \delta(C)=
 \dim_K H^0\bigl(C,\nu_*\cO_{\widetilde C}/\cO_C\bigr).
 \label{eq:conductor-delta}
\end{equation}
The normalization may be disconnected.  The comparison below follows the
evaluation-ideal argument of Fantechi--G\"ottsche--van Straten
\cite[Proposition C.2 and Corollary C.3]{FGvS1999}.
Their setting is over \(\C\); we give the local proof needed here in
arbitrary characteristic.

\begin{lemma}
\label{lem:conductor-local}
Let \(k\ge1\) be an integer, let \(K\) be an algebraically closed field,
and let \(S\) be a discrete
valuation ring containing \(K\), with residue field \(K\) and fraction
field \(Q\).  For \(A_0\in\Mat_k(S)\), put
\[
 M_0=yI+A_0,\qquad f=\det M_0,\qquad B=S[y]/(f).
\]
Suppose that \(f\) is separable over \(Q\).  Let \(R\) be the normalization
of \(B\) in \(B_Q=B\otimes_S Q\), let
\[
 \mathfrak c=\{a\in B:aR\subset B\},
\]
and let \(J_0\subset B\) be the ideal generated by the entries of
\(\adj M_0\).  Then
\begin{equation}
 \mathfrak c\subset J_0,\qquad
 \length_S(B/J_0)\le
 \length_S(B/\mathfrak c)=\length_S(R/B).
 \label{eq:conductor-local-length}
\end{equation}
\end{lemma}

\begin{proof}
Monicity makes \(B\) free over \(S\), with basis
\(1,y,\ldots,y^{k-1}\).  In particular, it embeds into \(B_Q\), a product
of finite separable field extensions of \(Q\), and is reduced.  The integral closures of
\(S\) and \(B\) in this algebra agree, since \(B\) is integral over \(S\).
They are finite over \(S\).  One way to see the finiteness directly is to
use the sum of the field trace pairings on the factors of \(B_Q\).
This pairing is nondegenerate, so the lattice dual to \(B\) is finite
free over \(S\).  If \(z\in B_Q\) is integral over \(S\), then each
\(zy^j\), for \(0\le j<k\), is integral because \(y\) is integral.
Its trace lies in the integrally closed ring \(S\), so \(z\) belongs
to the dual lattice.
The normalization is an \(S\)-submodule of that finite lattice and is
therefore finite.  The same finiteness conclusion follows by applying
the normalization theorem \cite[Tag 032L]{StacksProject} to each field
factor.

It follows that \(R\) is semilocal.  Its localizations at maximal ideals
are one-dimensional Noetherian normal local domains, hence discrete
valuation rings \cite[Tag 00PD]{StacksProject}.  All their residue fields
are \(K\).  As a finite torsion-free \(S\)-module, \(R\) is free of
rank \(k\), and \(R/B\) has finite length.

\medskip
\noindent
\emph{The adjugate ideal as an evaluation ideal.}
Consider the \(B\)-module
\[
 L=\operatorname{coker}(M_0:B^k\longrightarrow B^k).
\]
The same cokernel over \(S[y]\) is \(S^k\), with \(y\) acting as
\(-A_0\): polynomial-vector division by the monic matrix \(yI+A_0\)
reduces every class uniquely to a constant vector.  The adjugate identity
shows that \(f\) already kills this cokernel.  Thus \(L\) is free of
rank \(k\) over \(S\).

Over a splitting field of \(f\), the action of \(-A_0\) has \(k\)
distinct eigenvalues, each with a one-dimensional eigenspace.
Consequently \(L_Q\) is rank one over every field factor of \(B_Q\),
and \(L_Q\simeq B_Q\) as a \(B_Q\)-module.  Fixing such an isomorphism
embeds the \(S\)-free module \(L\) as a lattice in \(B_Q\).

Define its evaluation ideal by
\[
 \tau_B(L)=
 \operatorname{im}\left(
 \operatorname{Hom}_B(L,B)\otimes_B L\longrightarrow B
 \right).
\]
A homomorphism \(L\to B\) is a row annihilating \(M_0\) modulo \(f\).
Lift the row to \(b\in S[y]^{1\times k}\), so that
\(bM_0=fb'\) for another polynomial row \(b'\).
Multiplication by the adjugate gives
\[
 fb=fb'\adj M_0.
\]
We may cancel the nonzero \(f\) in the domain \(S[y]\), obtaining
\(b=b'\adj M_0\).  Conversely, each adjugate row annihilates \(M_0\)
modulo \(f\).  Evaluation on the standard generators of \(L\) therefore
gives
\begin{equation}
 \tau_B(L)=J_0.
 \label{eq:conductor-evaluation}
\end{equation}

\medskip
\noindent
\emph{The conductor is contained in the evaluation ideal.}
The fractional ideal \(I=RL\subset B_Q\) is free of rank one at every
maximal localization of \(R\).  We claim that some \(h\in L\) generates
it at all these localizations at once.  Choose \(S\)-generators
\(l_1,\ldots,l_k\) of \(L\).  At a maximal ideal \(\mathfrak m\) of \(R\),
their images span the one-dimensional \(K\)-space
\(I_{\mathfrak m}/\mathfrak m I_{\mathfrak m}\).  The coefficients
\((a_1,\ldots,a_k)\in K^k\) for which \(\sum_i a_i l_i\) has zero
image form a proper linear subspace.  There are only finitely many such
maximal ideals.  Since \(K\) is infinite, we can avoid all these subspaces.
The resulting \(h=\sum_i a_i l_i\in L\) generates every \(I_{\mathfrak m}\),
and hence \(I=hR\).

In particular, \(h\) is nonzero in every field factor of \(B_Q\).
The rescaled module \(L'=h^{-1}L\) satisfies
\[
 B\subset L'\subset R:
\]
it contains \(1\), and \(L\subset RL=hR\).
For \(a\in\mathfrak c\), multiplication by \(a\) defines a homomorphism
\(L'\to B\) whose value at \(1\) is \(a\).
Evaluation ideals are invariant under module isomorphism, so
\[
 \mathfrak c\subset\tau_B(L')
 =\tau_B(L)=J_0.
\]

We have identified the adjugate ideal with an evaluation ideal and shown
that it contains the conductor. It remains to compute the conductor
colength in terms of the normalization.

\medskip
\noindent
\emph{The conductor colength.}
Let \(\lambda:B\to S\) extract the coefficient of \(y^{k-1}\) in
the monic remainder, and consider
\begin{equation}
 \Phi:B\longrightarrow B^\vee:=\operatorname{Hom}_S(B,S),
 \qquad
 \Phi(a)(b)=\lambda(ab).
 \label{eq:conductor-pairing}
\end{equation}
This is \(B\)-linear for the action
\((a\cdot\ell)(b)=\ell(ab)\) on \(B^\vee\).
In the basis \(1,y,\ldots,y^{k-1}\), its matrix has entry
\(\lambda(y^{i+j})\) in position \((i,j)\).
That entry is zero when \(i+j<k-1\), and is one when \(i+j=k-1\).
Reversing the columns therefore gives a triangular matrix with diagonal
entries one.  Thus \(\Phi\) is an isomorphism in every characteristic.
The trace pairing above was used over \(Q\) to prove finiteness of \(R\).
This coefficient pairing supplies the self-duality over \(S\) used below.

The adjunction
\[
 \operatorname{Hom}_S(R,S)
 \simeq
 \operatorname{Hom}_B(R,\operatorname{Hom}_S(B,S))
\]
sends \(\ell\) to the map
\(x\mapsto[b\mapsto\ell(bx)]\).
Using \(\Phi\), it identifies \(\operatorname{Hom}_S(R,S)\) with
\(\operatorname{Hom}_B(R,B)\).
Every homomorphism in the latter module becomes a \(B_Q\)-linear
endomorphism of \(B_Q\), hence multiplication by some \(a\in B_Q\).
It takes \(R\) into \(B\) exactly when \(aR\subset B\).
Since \(1\in R\), this forces \(a\in B\), and the possible multipliers
are precisely \(\mathfrak c\).

This identification respects restriction to \(B\): for a multiplier \(a\),
the corresponding functional restricts to
\(b\mapsto\lambda(ab)=\Phi(a)(b)\).
Hence restriction
\begin{equation}
 R^\vee=\operatorname{Hom}_S(R,S)\longrightarrow B^\vee
 \quad\hbox{has image }\Phi(\mathfrak c).
 \label{eq:conductor-dual-image}
\end{equation}
To compute its colength, choose bases for the two free \(S\)-lattices
in which the inclusion \(B\subset R\) has elementary divisors
\(\pi^{n_1},\ldots,\pi^{n_k}\), where \(\pi\) is a uniformizer and
the \(n_i\) are nonnegative.  The quotient \(R/B\) has length
\(\sum_i n_i\), and the restriction on dual lattices has the same
elementary divisors.  By \eqref{eq:conductor-pairing} and
\eqref{eq:conductor-dual-image}, its cokernel is isomorphic to
\(B/\mathfrak c\).  This proves
\[
 \length_S(B/\mathfrak c)=\length_S(R/B).
\]
The inclusion \(\mathfrak c\subset J_0\) now gives
\eqref{eq:conductor-local-length}.
\end{proof}

\begin{proposition}
\label{prop:conductor-length}
For each curve \(C=C_X\) and adjugate subscheme \(Z\) constructed in
Section~\ref{sec:adjugate},
\begin{equation}
 \length Z\le\delta(C).
 \label{eq:conductor-global-length}
\end{equation}
\end{proposition}

\begin{proof}
By Lemma~\ref{lem:adjugate-projection}, the projection
\(\pi:C\to\PP^1_K\) is finite and flat, with a distinct-root fiber.
It is therefore generically separable.  At a closed point \(b\) of the
base, take \(S=\cO_{\PP^1,b}\) and trivialize \(\cO(1)\).
If \(y\) is the resulting \(s\)-coordinate and \(A_0\) represents
\(tD+uX\), the curve algebra is
\[
 B=S[y]/\bigl(\det(yI+A_0)\bigr).
\]
The stalk of \(\pi_*(\cO_C/J)\) is \(B/J_0\), with \(J_0\) as in
Lemma~\ref{lem:conductor-local}.  Its normalization \(R\) gives the
stalk \(R/B\) of
\(\pi_*(\nu_*\cO_{\widetilde C}/\cO_C)\).
Indeed, normalization commutes with this localization.
Changes of trivialization multiply the adjugate generators by units,
so these are the intrinsic ideals and quotient sheaves.

Both quotient sheaves have finite support.  At every closed point of
the base, their finite-length \(S\)-modules have length equal to their
\(K\)-dimension, because the residue field is \(K\).
Equivalently, each such length is the sum of the lengths at the curve
points above that base point; no ramification factor enters.
Summing \eqref{eq:conductor-local-length} over \(b\) therefore gives
\[
 \length Z
 =\sum_b\length_{\cO_{\PP^1,b}}(B/J_0)
 \le
 \sum_b\length_{\cO_{\PP^1,b}}(R/B)
 =\delta(C).
\]
Only generic separability was needed.  Nonreduced closed fibers and
wild ramification are allowed by the coefficient pairing
\eqref{eq:conductor-pairing}.
\end{proof}

In particular, Proposition~\ref{prop:adjugate-length} yields
\begin{equation}
 \nabla G_D(X)=0
 \quad\Longrightarrow\quad
 \delta(C_X)\ge T.
 \label{eq:conductor-gradient-defect}
\end{equation}
